\documentclass[10pt,aspectratio=169]{beamer}
\usepackage[UTF8]{ctex}
\usetheme{Berlin}
\usecolortheme{sustech}
\usepackage{amsmath,amssymb,booktabs,listings,xcolor}
\usepackage{hyperref}

\definecolor{codegray}{RGB}{245,245,245}
\lstset{
  language=C++,
  basicstyle=\ttfamily\scriptsize,
  backgroundcolor=\color{codegray},
  breaklines=true,
  columns=fullflexible
}

\newcommand{\imgplaceholder}[1]{%
  \fbox{\parbox[c][3.0cm][c]{0.92\linewidth}{\centering\small 占位图：#1}}
}

\title{AI算力中心任务调度问题题解}
\subtitle{在线优化建模 + 启发式最优化策略}
\author{讲题人：你}
\institute{中关村文理学院（模拟场景）}
\date{20分钟讲解版}

\begin{document}

\begin{frame}
  \titlepage
\end{frame}

\begin{frame}{目录}
  \tableofcontents
\end{frame}

\section{题意与建模}

\begin{frame}{1. 背景与建模对象（1.5 min）}
\textbf{场景：}校内 AI 计算中心给同学提供免费算力，调度程序要在主机上实时响应请求。\\[4pt]
\textbf{系统总资源：}
\begin{itemize}
  \item 最多同时运行 $n$ 个任务（进程位）
  \item 总 GPU 数 $K$（同质显卡）
  \item 总共享内存 $M$ GB
\end{itemize}
\textbf{任务公开信息：}
\begin{itemize}
  \item 需求 $(k_i,m_i)$、段长 $t_i$、重要性 $W_i$
  \item 任务编号递增，进程编号 $0\sim n-1$
\end{itemize}
\imgplaceholder{“任务流 + 多资源池 + 进程槽位”框图}
\end{frame}

\begin{frame}{2. 任务模型与目标函数（2 min）}
任务 $i$ 的 checkpoint 机制：
\begin{itemize}
  \item 开始时刻算第一个 checkpoint
  \item 相邻 checkpoint 间隔恒为 $t_i$（整小时）
  \item checkpoint 总数 $c_i$ \textbf{开始未知}，仅在任务结束时获知
  \item 题目简化下总时长严格为 $c_i t_i$
\end{itemize}
目标函数：
\[
\min \sum_i \frac{W_i T_i}{c_i t_i}
\]
其中 $T_i$ 是从任务发出到最终完成的总周转时间。
\end{frame}

\begin{frame}{3. checkpoint 与抢占语义（1.5 min）}
\begin{itemize}
  \item 允许抢占，抢占本身不消耗时间。
  \item 但若不在 checkpoint 边界抢占，会丢失“上个 checkpoint 后”全部进度。
  \item 所以段内抢占的真实代价与“已跑段长”正相关。
\end{itemize}
\vspace{4pt}
调度核心矛盾：
\begin{itemize}
  \item 继续跑旧任务：保护已投入进度
  \item 抢占上新任务：降低高价值任务等待损失
\end{itemize}
\imgplaceholder{“段内进度回滚 + checkpoint 无损切换”时间轴}
\end{frame}

\begin{frame}{4. 交互机制、合法性与评分（2 min）}
\textbf{接口：}
\texttt{Init(...)}、\texttt{Event(now,event)}，并可调用 \texttt{Schedule} / \texttt{Appoint}。\\[2pt]
\textbf{事件：}checkpoint / 结束 / 新任务 / 预约事件。\\[2pt]
\textbf{合法性红线：}
\begin{itemize}
  \item \texttt{Schedule} 时任务不能已在别处运行，且资源必须可满足
  \item \texttt{Appoint(appoint\_time)} 要满足 $appoint\_time\ge now$
\end{itemize}
\textbf{评分：}
\[
score=15\cdot\min\!\left(\frac{0.95\cdot std\_perf}{your\_perf},1\right)
\]
每点时限 8s（C++），在线算法必须兼顾效果与开销。
\end{frame}

\section{题解推导与策略设计}

\begin{frame}{5. 从目标函数推导“单位等待边际代价”（1.5 min）}
先看单任务代价项：
\[
F_i=\frac{W_i T_i}{c_i t_i}
\]
若任务 $i$ 额外等待 $\Delta t=1$ 小时，则近似增量：
\[
\Delta F_i \approx \frac{W_i}{c_i t_i}
\]
这就是任务 $i$ 当前“再等 1 小时”的边际损失。\\[4pt]
因此在线调度直觉：
\begin{block}{自然优先级原则}
优先运行边际损失 $\frac{W_i}{c_i t_i}$ 更大的任务。
\end{block}
\end{frame}

\begin{frame}{6. 关键难点：$c_i$ 未知，改成在线估计（1.5 min）}
由于真实 $c_i$ 在完成前不可见，定义估计：
\[
\hat c_i=\text{checkpoints\_passed}_i+\text{mean\_segments}
\]
得到可计算的价值密度：
\[
\phi_i=\frac{W_i}{\hat c_i t_i}
\]
\begin{itemize}
  \item $\phi_i$ 越大，说明“让它再等 1 小时”预计越亏
  \item 新任务用全局先验 \texttt{mean\_segments}
  \item 已跑过多段却未完成的任务，$\hat c_i$ 会自然变大
\end{itemize}
\end{frame}

\begin{frame}{7. 全局先验如何更新：指数平滑（1.5 min）}
任务完成时可观测到其总段数近似：
\[
c_i\approx \text{checkpoints\_passed}_i+1
\]
用 EMA 更新全局先验：
\[
\text{mean\_segments}\leftarrow(1-\alpha)\cdot\text{mean\_segments}
+\alpha\cdot(\text{checkpoints\_passed}_i+1)
\]
\begin{itemize}
  \item $\alpha$ 小：更平滑、更稳
  \item $\alpha$ 大：更快跟踪分布漂移
\end{itemize}
\textbf{对应代码参数：} \texttt{MEAN\_SEGMENTS\_PRIOR} + \texttt{MEAN\_SEGMENTS\_ALPHA}。
\end{frame}

\begin{frame}{8. 当前时刻的理想决策：约束最优化（2 min）}
选运行集合 $S$，满足：
\[
|S|\le n,\quad \sum_{i\in S}k_i\le K,\quad \sum_{i\in S}m_i\le M
\]
当前未运行任务每小时损失约为 $\phi_i$，故瞬时增长率：
\[
\sum_{i\notin S}\phi_i=\sum_i\phi_i-\sum_{i\in S}\phi_i
\]
前项为常数，等价为：
\[
\max_{S}\ \sum_{i\in S}\phi_i
\]
\textbf{结论：}这是一个三重约束（数量/GPU/内存）的在线装箱-背包问题。
\end{frame}

\begin{frame}{9. 启发式贪心：价值与资源体积的比值（2 min）}
精确求最优过慢，采用可在线重算的贪心评分。\\[-2pt]
定义资源体积：
\[
d_i=\alpha\frac{k_i}{K}+(1-\alpha)\frac{m_i}{M}
\]
定义排序分数：
\[
score_i=\frac{\phi_i}{d_i\cdot t_i^{\theta}}
\]
\begin{itemize}
  \item 分子 $\phi_i$：等待边际损失
  \item 分母 $d_i$：资源占用强度
  \item $t_i^{\theta}$：轻微长段惩罚，提高调度灵活性
\end{itemize}
按 $score$ 降序尝试装入，不违反 $(n,K,M)$ 就选入。
\end{frame}

\begin{frame}{10. 为什么要加 $t_i^{\theta}$ 惩罚（1 min）}
若 $t_i$ 很大：
\begin{itemize}
  \item 下一次无损切换点（checkpoint）来得更慢
  \item 调度器更容易被“长段任务”卡住
\end{itemize}
因此加轻惩罚（如 $\theta=0.2$）：
\[
score_i\propto \frac{1}{t_i^{\theta}}
\]
不是追求公平，而是控制段内抢占风险、提升全局可调度性。
\end{frame}

\begin{frame}{11. 从理想集合 $S$ 到真实进程映射（1.5 min）}
\textbf{Step 1：}保留“当前在跑且仍属于 $S$”的任务（不动最省）。\\
\textbf{Step 2：}构造 \texttt{need\_start}：属于 $S$ 但未运行的任务。\\
\textbf{Step 3：}逐个安置 \texttt{need\_start}：
\begin{itemize}
  \item 优先放空闲进程（无抢占成本）
  \item 放不下再考虑抢占不在 $S$ 中的运行任务
\end{itemize}
\textbf{工程意义：}把“理论最优集合”变成“最小扰动的可执行调度”。
\end{frame}

\begin{frame}{12. 段内抢占阈值：收益必须覆盖损失（2 min）}
新任务 $i$ 想抢占运行任务 $j$，记：
\[
s_j=\text{本段已运行时长},\quad r_j=t_j-s_j
\]
若立即抢占：损失约 $\phi_j s_j$。\\
若等到下个 checkpoint 再切：新任务多等 $r_j$，损失约 $\phi_i r_j$。\\[2pt]
仅当
\[
\phi_i r_j > \lambda\,\phi_j s_j,\quad \lambda\ge1
\]
才允许段内抢占。\\
\textbf{解释：}新任务足够“贵”且旧任务已跑段损失可接受，才值得打断。
\end{frame}

\begin{frame}{13. 多候选可抢占时：选最小驱逐代价（1.5 min）}
对候选被驱逐任务 $j$ 定义：
\[
\text{evict\_cost}(j)=\lambda_1\phi_j s_j+\lambda_2\phi_j r_j
\]
\begin{itemize}
  \item 第一项：已跑段回滚损失
  \item 第二项：临近 checkpoint 的保护项
\end{itemize}
在“资源可行 + 通过抢占阈值”的候选进程中，取 \texttt{evict\_cost} 最小者。
\imgplaceholder{“多个候选进程的驱逐代价柱状对比图”}
\end{frame}

\begin{frame}{14. 事件驱动状态维护（1.5 min）}
\textbf{type=3 新任务：}创建 WAITING，记录 \texttt{waiting\_since=now}。\\
\textbf{type=1 checkpoint：}\texttt{checkpoints\_passed++}，刷新段起点。\\
\textbf{type=2 完成：}释放资源、更新 \texttt{mean\_segments}、删除任务。\\
\textbf{type=4 预约：}仅触发一次重调度（本策略不强依赖）。\\[4pt]
\textbf{等待时间维护：}
\[
T_i=\begin{cases}
waiting\_accum,& RUNNING\\
waiting\_accum+(now-waiting\_since),& WAITING
\end{cases}
\]
每次状态切换可 $O(1)$ 更新。
\end{frame}

\begin{frame}{15. 复杂度与时限把控（1.5 min）}
设当前未完成任务数为 $N$：
\begin{itemize}
  \item 计算评分：$O(N)$
  \item 选前列候选 + 局部排序：近似 $O(N\log N)$（可用 \texttt{nth\_element} 降常数）
  \item 贪心装箱：$O(N)$
  \item 进程映射/候选抢占：与进程数 $n$ 线性相关
\end{itemize}
单次事件主复杂度约 $O(N\log N)$，可覆盖题目 1e4 任务规模。\\
工程上还可采用“并非每次事件都全量重算”的降频策略。
\end{frame}

\begin{frame}{16. 总结：最终策略的一句话（1 min）}
\begin{block}{核心结论}
每次事件时重估“谁再等 1 小时最亏”（$\phi_i$），在 $(n,K,M)$ 约束下选最值得跑的集合；已在跑且仍应跑的尽量不动；仅当新任务收益显著大于丢段代价时才段内抢占，并选择最小驱逐代价目标。
\end{block}
\vspace{4pt}
\centering
\Large Thanks \& Q\&A
\end{frame}

\end{document}
